\subsection{Installation}

\subsubsection{Python environment and IDEs}

The Python distribution \textbf{Anaconda} is recommended. To install \textbf{Anaconda}, see \href{https://www.anaconda.com/docs/getting-started/anaconda/install/overview}{Installing Anaconda Distribution}. After installation, a \textbf{conda} virtual environment can be created by, e.g.,

\begin{verbatim}
conda create - -name YOUR_ENV_NAME python=3.12 cython spyder
\end{verbatim}

After creating the virtual environment, one needs to activate it:

\begin{verbatim}
conda activate YOUR_ENV_NAME
\end{verbatim}

\paragraph{Python IDEs}

Popular Python IDEs include \textbf{Spyder}, \textbf{PyCharm}, and \textbf{VS Code}.

\subsubsection{Install GeospaceLAB}

Install GeospaceLAB via PyPI:

\begin{verbatim}
pip install geospacelab
\end{verbatim}

To upgrade GeospaceLAB:

\begin{verbatim}
pip install -U geospacelab
\end{verbatim}

\paragraph{Dependencies}

Most of the dependencies are installed when the package is installed. A few packages can be installed manually.

For example, the \textbf{Cartopy} package is useful for processing and mapping geospatial data. To install \textbf{Cartopy}, see \href{https://cartopy.readthedocs.io/stable/installing.html}{Installing Cartopy}.

Another useful package is \textbf{ApexPy}, which provides a Python interface to the Apex coordinate system. To install \textbf{ApexPy}, see \href{https://apexpy.readthedocs.io/en/latest/installation.html}{Installing ApexPy}.

\begin{framed}
\textbf{Note}\\
To use the VirES and HAPI services, install the \texttt{viresclient}:

\begin{verbatim}
pip install viresclient
\end{verbatim}

and the \texttt{hapiclient}:

\begin{verbatim}
pip install hapiclient
\end{verbatim}
\end{framed}